\documentclass[10pt,twocolumn]{article}
\usepackage[margin=0.75in]{geometry}
\usepackage{amsmath,amssymb}
\usepackage{booktabs}
\usepackage{listings}
\usepackage{xcolor}
\usepackage{hyperref}
\hypersetup{colorlinks=true,linkcolor=blue,citecolor=blue,urlcolor=blue}

\lstset{
  basicstyle=\ttfamily\footnotesize,
  breaklines=true,
  frame=single,
  framesep=3pt,
  columns=fullflexible,
  keywordstyle=\color{blue!60!black},
  commentstyle=\color{gray},
  language=Go,
}

\title{\textbf{Measured Primitives}\\[2pt]
\large Three Zero-Dependency Go Packages for Forecasting,\\Memoization and Deferred Evaluation}

\author{Daniel Paes\\\small Runink\\\small\texttt{github.com/org-runink/runi}}
\date{\today}

\begin{document}
\maketitle

\begin{abstract}
We present \texttt{stdx}, three small Go packages built under a single
constraint: no dependencies outside the standard library. The constraint is not
aesthetic. It is what makes the packages usable in a CLI, a sidecar, a WASM
target or a device build, and it forces every numerical and concurrency decision
to be made explicitly rather than delegated.

We make three claims and support each with a reproducible measurement. First,
that an ARIMAX implementation should \emph{refuse} to forecast without the future
values of its exogenous regressors, rather than silently assuming zero or holding
the last value flat as implementations commonly do; and that its prediction
intervals should be validated empirically rather than merely derived---ours
achieve $97.4\%$ coverage at a nominal $95\%$ level over $1{,}800$ held-out
points. Second, that single-flight coalescing belongs \emph{inside} a memoization
cache rather than beside it, because sequential repeats and concurrent duplicates
are disjoint populations: our hit path costs $14.34$\,ns with zero allocations,
and $64$ concurrent callers on a cold key cost one execution. Third, that lazy
evaluation should separate \emph{whether} work happens from \emph{when} it
starts; speculative start reduces a forced-value latency from $122$\,ms to
$8\,\mu$s without sacrificing the property that an unforced value is never
computed.

We also report what we did not measure, and one negative result: in our own
deployment, redundant computation was reduced substantially more by declining to
invoke an expensive model than by replaying its prior answers, which suggests
memoization and avoidance are complementary rather than competing strategies.
\end{abstract}

\section{Introduction}

Three recurring costs motivate this work: forecasting a quantity that is driven
by external factors, repeating a computation whose inputs have not changed, and
paying latency for a value that could have been computed earlier. Each has a
textbook answer. In practice each textbook answer is implemented in a way that
fails quietly.

An ARIMAX implementation that is asked to forecast without future regressor
values will often proceed anyway. A cache helps the second caller while the
first $N$ arrive together and each pays full cost. A lazy value defers work
until it is forced, and then the forcing caller waits for all of it.

The packages described here are small---under $1{,}500$ lines of
non-test Go in total---and the contribution is not algorithmic novelty. Every
method used predates 1980. The contribution is that each failure mode above is
treated as a design requirement and each claim is attached to a test that fails
when the claim stops being true.

\section{Design constraint: no dependencies}

\texttt{go.mod} declares no requirements. This is verified in CI by
\texttt{go list -deps}.

The cost is that Householder QR, Nelder--Mead, the Durbin--Levinson recursion
and an inverse normal CDF are implemented rather than imported; each is a few
dozen lines and each is covered by a known-answer test. The benefit is
deployability and, less obviously, that every numerical choice had to be made and
can therefore be documented.

\section{\texttt{arimax}: forecasting with external drivers}

ARIMA models a series by its own past. ARIMAX adds exogenous regressors:
\begin{equation}
Y_t = c + \sum_{i=1}^{p}\phi_i Y_{t-i} + \sum_{j=1}^{q}\theta_j \varepsilon_{t-j}
      + \sum_{k=1}^{r}\beta_k X_{t,k} + \varepsilon_t
\end{equation}

\subsection{Staged estimation}

We fit by the regression-with-ARIMA-errors formulation rather than a joint
likelihood: ordinary least squares for $\beta$, then an ARMA fit to the
regression errors $n_t = Y_t - \beta'X_t - c$. The two are equivalent in what they
express; the staged form is better conditioned and each stage is separately
testable.

Least squares uses Householder QR rather than the normal equations. Exogenous
regressors in applied settings are routinely near-collinear---advertising spend
across channels, say---and forming $X'X$ squares the condition number.

ARMA parameters minimise the conditional sum of squares by Nelder--Mead. CSS
conditions on the first $p$ observations rather than running a Kalman filter for
the exact likelihood. We state this rather than implying maximum likelihood.

\subsection{Forecasting requires future $X$}

A forecast $h$ steps ahead requires $X_{t+1},\dots,X_{t+h}$. Our API takes them
as an argument and returns an error when they are absent.

\begin{lstlisting}
point, lo, hi, err := m.Forecast(12, xFuture, 0.05)
// err if xFuture is nil or mis-sized
\end{lstlisting}

Where $X$ is known in advance---committed spend, a published timetable---this is
the model at its best. Where $X$ is itself forecast, its error propagates and the
intervals below \emph{understate} the true uncertainty because they condition on
$X$ being correct. We document this rather than leaving it to be discovered.

\subsection{Results}

Parameters are recovered from $200$ independent synthetic series with AR(1)
errors and one regressor, $n=500$. Synthetic data is used deliberately: bias
cannot be measured without knowing the truth.

\begin{table}[h]\centering\small
\begin{tabular}{lrrr}
\toprule
Parameter & true & bias & RMSE \\
\midrule
exogenous $\beta$ & 2.50 & $-0.0005$ & 0.0087 \\
AR(1) $\phi$      & 0.60 & $-0.0079$ & 0.0420 \\
\bottomrule
\end{tabular}
\caption{Parameter recovery. $\beta$ is recovered essentially unbiased despite
strongly serially correlated errors, which is the purpose of the staged fit. The
small negative bias in $\phi$ is the known finite-sample bias of CSS estimation
and is not corrected.}
\end{table}

\noindent\textbf{Interval coverage.} Over $300$ series and $1{,}800$ held-out
forecast points at horizons $1..6$, nominal $95\%$ prediction intervals contained
the realised value $97.4\%$ of the time.

An interval that is never validated is decoration. One that claims $95\%$ and
delivers $70\%$ is worse than none, because it invites confident wrong answers.
Our intervals are mildly \emph{conservative}, which is the safe direction; the
gap arises because they assume known parameters and omit estimation error, while
the CSS residual variance is slightly inflated in finite samples.

\noindent\textbf{Against a naive baseline.} Over $100$ held-out windows at
$h{=}6$, mean RMSE was $0.6188$ versus $3.4424$ for last-value carry-forward,
$82.0\%$ lower. We report this as a floor rather than an achievement: a
forecaster that cannot beat ``tomorrow equals today'' is not earning its
complexity.

\section{\texttt{memo}: memoization with single-flight}

\subsection{Why coalescing belongs inside the cache}

A cache serves the second caller. The characteristic production failure is
different: $N$ callers arrive concurrently on a cold key, all miss, and all
execute. Sequential repeats and concurrent duplicates are disjoint populations
and neither mechanism covers the other, so we implement both in one type.

\subsection{Errors are coalesced but not stored}

A failed call is shared with everyone waiting on it, and is \emph{not} cached.
Caching a failure converts a transient fault into a sticky one for the whole TTL;
in our experience this is the most common way a memoization cache makes a system
worse rather than better. A stampede against a failing dependency still produces
one call rather than $N$.

\subsection{Key canonicalisation}

The difficult part of exact-match caching is deciding that two requests are the
same request. Go randomises map iteration order, so hashing a map naively yields
a different key on every run---a cache that never hits and never reports why. Our
\texttt{Hash} sorts map keys, type-tags every value so that \texttt{"1"}, \texttt{1}
and \texttt{true} differ, length-prefixes variable-length data so that
$\{\texttt{"ab"},\texttt{"c"}\}$ and $\{\texttt{"a"},\texttt{"bc"}\}$ cannot
collide, and canonicalises NaN and $-0.0$.

\subsection{Results}

\begin{table}[h]\centering\small
\begin{tabular}{lrrr}
\toprule
Operation & ns/op & B/op & allocs \\
\midrule
hit                       & 14.34 & 0 & 0 \\
miss (store + LRU insert) & 694.2 & 358 & 5 \\
hit, 16 threads           & 134.9 & 9 & 0 \\
\texttt{Hash}, 4-field map & 1579 & 656 & 36 \\
stampede, 64 cold callers & 24{,}324 & 4136 & 75 \\
\bottomrule
\end{tabular}
\caption{The hit path allocates nothing. The parallel figure is $9\times$ the
single-threaded one: \texttt{Store} takes one mutex per operation. We report this
rather than omit it---at roughly 7M hits/s aggregate it is irrelevant for network
or model calls, and disqualifying for sub-microsecond work.}
\end{table}

\section{\texttt{lazy}: deferred values that can start early}

Laziness decides \emph{whether} work happens. Starting decides \emph{when}. Most
lazy types conflate the two, so the caller who forces a value pays its full cost.
We separate them: \texttt{Start} begins evaluation in the background without
forcing, and an unforced value is still never computed.

\begin{lstlisting}
v := lazy.New(build)
v.Start(ctx)         // returns immediately
...                  // other work
r, err := v.Get(ctx) // already resolved
\end{lstlisting}

\begin{table}[h]\centering\small
\begin{tabular}{lr}
\toprule
Measurement & \\
\midrule
\texttt{Get} on a cold value & 122\,ms \\
\texttt{Get} after \texttt{Start} completed & $8\,\mu$s \\
\texttt{All} over five 80\,ms values & 81\,ms \\
\quad (serial equivalent) & 400\,ms \\
\texttt{Get} on a resolved value & 73.91\,ns \\
\bottomrule
\end{tabular}
\caption{Speculative start removes the forcing caller's wait almost entirely.}
\end{table}

A panic inside the function is recovered and returned to every caller as an error
wrapping \texttt{ErrPanic}. Which goroutine happens to be forcing a shared value
is a scheduling accident, so allowing a panic to propagate would make the crash
site arbitrary.

\section{A negative result}

We audited our own production deployment for memoization and found none. What the
system had instead was avoidance: single-flight coalescing at admission, durable
idempotency for mutating replays, and a cascade of calibrated tree-based decisions
that defer to an expensive model only when uncertain.

Redundant computation was reduced substantially more by \emph{not invoking} the
model than by replaying its answers. This is not an argument against memoization.
It is evidence that the two are complementary and cover different cases:
coalescing addresses concurrent duplicates, cascades address cheaply-answerable
questions, and memoization addresses the sequential repeat that neither reaches.
A system optimised only for replay leaves the first two on the table, and one
optimised only for avoidance leaves the third.

\section{Threats to validity}

\textbf{No comparison against other libraries.} We report only our own
measurements. No claim of superiority over any alternative is made, because no
such experiment was run.

\textbf{Accuracy is measured on synthetic data.} This is the only way to measure
estimator bias, since it requires knowing the truth. It is not evidence of
accuracy on real series, where the model is misspecified by construction.

\textbf{Single machine.} All figures are from one developer laptop
(AMD Ryzen 7 8840U, Linux, \texttt{amd64}) with no core pinning. They indicate
shape---linear growth, allocation counts---not capacity.

\textbf{CSS is not maximum likelihood.} The two agree closely at the series
lengths benchmarked, but they are different estimators.

\textbf{Gaussian errors are assumed} by the prediction intervals; heavy-tailed
residuals will produce intervals that are too narrow in the tails.

\section{Related work}

Box and Jenkins~\cite{boxjenkins} established the ARIMA family; Hyndman and
Athanasopoulos~\cite{hyndman} give the practical treatment of regression with
ARIMA errors that we follow. Robertson et al.~\cite{bm25} introduced BM25, which
we use elsewhere in our retrieval stack but not in these packages. Recent work on
agentic search argues that unbounded corpus interaction does not scale and that
bounded retrieval remains necessary~\cite{rise}; we note it because it bears on
how systems built on these primitives should retrieve, not because these packages
implement it.

Memoization and single-flight are classical. Our contribution is their
combination with explicit error semantics and canonical key derivation, together
with the measurements above.

\section{Availability}

\texttt{github.com/org-runink/runi}, BSD-3-Clause. Every figure reproduces with
\texttt{go test -bench . -benchmem}. Statement coverage is $87.0\%$
(\texttt{arimax}), $82.4\%$ (\texttt{memo}) and $93.1\%$ (\texttt{lazy}).

\begin{thebibliography}{9}
\bibitem{boxjenkins} G.~E.~P. Box and G.~M. Jenkins. \emph{Time Series Analysis:
Forecasting and Control}. Holden-Day, 1970.
\bibitem{hyndman} R.~J. Hyndman and G. Athanasopoulos. \emph{Forecasting:
Principles and Practice}, 2nd ed. OTexts, 2018.
\bibitem{bm25} S.~E. Robertson et al. Okapi at TREC-3. \emph{TREC}, 1994.
\bibitem{rise} Zhuang, Ni, Fun, Lin and Ma. Towards Retrieving Interaction Spaces
for Agentic Search. arXiv:2606.06880, 2026.
\end{thebibliography}

\end{document}
